\documentclass[11pt,a4paper]{article}

\usepackage[utf8]{inputenc}
\usepackage[T1]{fontenc}
\usepackage[french]{babel}
\usepackage{geometry}
\usepackage{array}
\usepackage{longtable}
\usepackage[table]{xcolor}
\usepackage{hyperref}

\geometry{margin=2.2cm}

\definecolor{straussblue}{RGB}{36,63,196}
\definecolor{lightblue}{RGB}{234,245,255}

\title{\textbf{Scénarios utilisateur -- ARIA ODM Builder}}
\author{Projet ARIA ODM Builder}
\date{\today}

\begin{document}

\maketitle

\section{Objectif du document}

Ce document décrit les principaux scénarios d'utilisation de l'application
\textbf{ARIA ODM Builder}.  
L'objectif est de relier chaque besoin utilisateur à une action concrète dans
l'interface Streamlit.

Ces scénarios permettent de vérifier que la navigation de l'application est
fluide, que les onglets sont dans un ordre logique, et que les étapes importantes
ne sont pas cachées dans la barre latérale.

\section{Principe général}

L'application suit le parcours suivant :

\begin{center}
\textbf{Accueil} $\rightarrow$
\textbf{Import} $\rightarrow$
\textbf{Construction} $\rightarrow$
\textbf{Contrôle qualité} $\rightarrow$
\textbf{Sources} $\rightarrow$
\textbf{Profil}
\end{center}

Chaque onglet correspond à une étape du travail utilisateur.

\section{Scénarios utilisateur}

\renewcommand{\arraystretch}{1.35}

\begin{longtable}{|p{0.45\textwidth}|p{0.45\textwidth}|}
\hline
\rowcolor{lightblue}
\textbf{Narration utilisateur} & \textbf{Action correspondante dans l'application} \\
\hline
\endfirsthead

\hline
\rowcolor{lightblue}
\textbf{Narration utilisateur} & \textbf{Action correspondante dans l'application} \\
\hline
\endhead

L'utilisateur arrive sur l'application et veut comprendre l'objectif général de l'outil.
&
Il consulte l'onglet \textbf{Accueil}, qui présente le rôle de l'application, les fichiers attendus et le déroulé global.
\\
\hline

L'utilisateur doit importer les données brutes issues d'ARIA/MOSAIQ.
&
Il va dans l'onglet \textbf{Import} et charge le fichier \texttt{traitement\_patient}.
\\
\hline

L'utilisateur doit importer les données de formulaires patients.
&
Il reste dans l'onglet \textbf{Import} et charge le fichier \texttt{formulaire\_patient}.
\\
\hline

L'utilisateur dispose d'un fichier de correspondance entre les codes CIM10, les localisations et les colonnes utiles.
&
Il importe un fichier \texttt{mapping.csv}, idéalement placé dans le dossier \texttt{conf/}.
\\
\hline

L'utilisateur veut réutiliser une configuration déjà préparée lors d'une précédente session.
&
Il importe un profil JSON depuis l'onglet \textbf{Import} ou \textbf{Profil}.
\\
\hline

L'utilisateur souhaite construire une cohorte à partir d'un ou plusieurs codes CIM10.
&
Il va dans l'onglet \textbf{Construction}, choisit le mode CIM10, saisit les codes souhaités, puis clique sur \textbf{Construire / recalculer}.
\\
\hline

L'utilisateur veut éviter que l'application recalcule tout à chaque modification d'un paramètre mineur.
&
L'application attend une action explicite sur le bouton \textbf{Construire / recalculer} avant de relancer les traitements lourds.
\\
\hline

L'utilisateur souhaite sélectionner les colonnes formulaire pertinentes pour son export.
&
Dans l'onglet \textbf{Construction}, il utilise les suggestions issues du mapping et ajuste la sélection des colonnes.
\\
\hline

L'utilisateur veut définir les temporalités médicales associées aux variables sélectionnées.
&
Il coche les blocs temporels nécessaires : \textbf{Cumul}, \textbf{Avant RT}, \textbf{Aigu} et \textbf{Tardif}.
\\
\hline

L'utilisateur veut vérifier que les patients du fichier traitement sont bien retrouvés dans le fichier formulaire.
&
Il consulte l'onglet \textbf{Contrôle qualité}, qui affiche les contrôles de jointure, les dates manquantes et les doublons.
\\
\hline

L'utilisateur veut comprendre d'où viennent les colonnes générées dans l'export final.
&
Il consulte l'onglet \textbf{Sources}, qui documente les colonnes d'origine, les transformations et le schéma temporel généré.
\\
\hline

L'utilisateur veut exporter un fichier patient final exploitable.
&
Il utilise les boutons d'export pour générer le fichier final en CSV ou Excel.
\\
\hline

L'utilisateur veut conserver la preuve des choix effectués pendant la construction.
&
Il exporte un rapport de preuve contenant les paramètres, la qualité, le mapping et le schéma des colonnes générées.
\\
\hline

L'utilisateur veut sauvegarder sa configuration pour la réutiliser plus tard.
&
Il va dans l'onglet \textbf{Profil} et télécharge un profil JSON.
\\
\hline

L'utilisateur revient plus tard avec de nouveaux fichiers patients mais souhaite garder la même logique d'extraction.
&
Il importe les nouveaux fichiers, recharge le profil JSON, puis clique sur \textbf{Construire / recalculer}.
\\
\hline

\end{longtable}

\section{Conséquences sur la structure de l'interface}

Les scénarios précédents justifient les choix suivants :

\begin{itemize}
    \item l'onglet \textbf{Accueil} sert à expliquer le contexte et le flux général ;
    \item l'onglet \textbf{Import} est une étape dédiée, car l'import des fichiers est obligatoire avant toute construction ;
    \item l'onglet \textbf{Construction} contient les paramètres de cohorte, de sélection des colonnes et de temporalité ;
    \item l'onglet \textbf{Contrôle qualité} permet de vérifier la cohérence des données avant export ;
    \item l'onglet \textbf{Sources} rend l'export traçable ;
    \item l'onglet \textbf{Profil} permet de sauvegarder et réutiliser une configuration.
\end{itemize}

\section{Intérêt des scénarios}

Les scénarios permettent de passer d'une simple liste de fonctionnalités à une
description concrète du parcours utilisateur.

Ils servent à :

\begin{itemize}
    \item vérifier que l'ordre des onglets est logique ;
    \item éviter de cacher une étape obligatoire dans la barre latérale ;
    \item identifier les endroits où un bouton explicite est nécessaire ;
    \item améliorer la lisibilité de l'application pour un utilisateur non développeur ;
    \item faciliter les discussions avec les utilisateurs finaux et les encadrants.
\end{itemize}

\end{document}
